\section{Future work}
\label{sec:future}

Two natural extensions sit immediately downstream of the present
formalization, both prepared for by the v0.4 two-tier architecture
of \S\ref{sec:proof}.

\subsection{Geometric lift to a sweepout}
\label{sec:future-lift}

The structured output of \S\ref{sec:proof}
(\lean{exists\_DLTCover}{OneDim/Assembly.lean}) packages the Stage A
initial cover with spacing condition~(a), the Stage B partial
refinement, $\sigma$-injectivity, and the termination invariant
$\forall i,\ I_i \subseteq \bigcup_k J_k$ as a single Lean structure
\texttt{DLTCover}. The companion lift
\[
  \texttt{DLTCover}\;\longrightarrow\;
  \{\partial \Omega'_t\}_{t \in [0,1]}\ \text{flat-continuous}
\]
realizes the De~Lellis--Tasnady modified sweepout
$\Omega'_t$ defined by cut-paste blending on the positive-measure
overlaps $J_i \cap J_{i+1}$. The blending formula
\[
  \Omega'_t \;=\; [\Omega_t \setminus (U_i' \cup U_{i+1}')]
   \;\cup\; [\Omega_{i,t} \cap U_i']
   \;\cup\; [\Omega_{i+1,t} \cap U_{i+1}']
\]
is purely set-algebraic, and its $t$-continuity in the flat metric
follows from $t$-continuity of the input families
$\{\Omega_t\}, \{\Omega_{i,t}\}$ together with $t$-independence of
the cut sets $U_i', U_{i+1}'$; no smooth structure or partition of
unity is required. The lift therefore becomes formally available
once an ambient Lean library supplies a finite-perimeter / sweepout
API with a flat-continuity notion --- a strictly geometric
prerequisite, separate from the combinatorial argument formalized
here.

The placeholder directory \texttt{CombArg/Geometric/} marks the
import position the future lift module will occupy.
Remark~\ref{rem:why-construction} is the design contract:
\texttt{DLTCover}'s preserved spacing/overlap structure, not the
abstract scalar bound of Corollary~\ref{thm:main}, is what the
geometric lift consumes; the cheap partition route
(\texttt{CombArg.Scalar.exists\_refinement\_partition}) shipped as
its companion verifies by dependency-graph fact that the DLT
structure is essential exactly at this lift step and at no earlier
one.

\subsection{Multi-parameter sweepouts}
\label{sec:future-multiparam}

The cover construction in \S\ref{sec:proof} is specialized to the
$1$-parameter case $K = [0,1]$. Multi-parameter min-max
constructions ($K = [0,1]^m$ for $m \ge 2$, in particular Almgren
cycles in Marques--Neves~\cite{MN14}) replace the skip-$2$ spacing
of \S\ref{sec:proof} (an inequality on $1$-dimensional indices)
with a corresponding spacing structure on an $m$-dimensional grid;
the parity rescue (\S\ref{sec:proof}, Proposition~\ref{prop:twofold})
that delivers the two-fold overlap invariant from skip-$2$ spacing
needs an $m$-dimensional analogue. The bookkeeping corollary
\lean{exists\_sup\_reduction\_of\_cover}{Cover.lean} is already
$K$-generic in compact $K$, so only the cover construction itself
must be redone for the multi-parameter setting.

The natural development is a sister namespace
\texttt{CombArg.MultiDim} parallel to the present
\texttt{CombArg.OneDim}, producing a \texttt{DLTCover}-analogous
structured output for each $m \ge 1$ that the same bookkeeping
corollary then turns into a sup-reducing competitor on $K =
[0,1]^m$. The two-tier architecture inherited from the
$1$-dimensional case --- a structured DLT-style cover preserving
geometric data and a cheaper alternative scalar proof for the same
abstract conclusion --- is expected to carry over: in each
dimension, the abstract scalar reduction admits a partition-style
shortcut, while the geometric lift to a multi-parameter modified
sweepout consumes the structured form.

\medskip

Together, the two directions complete the formalization arc
sketched at the start of \S\ref{sec:integration}: this paper
formalizes the combinatorial slice; \S\ref{sec:future-lift} crosses
into geometric measure theory once an ambient library is in place;
and \S\ref{sec:future-multiparam} extends the combinatorial slice
itself across parameter dimensions. Each is independently scoped and
does not block the other.
